
% ----------------------------------------------------------------------
\section{Application : problème de régression}
% ----------------------------------------------------------------------

% ----------------------------------------------------------------------
\begin{frame}<beamer>
  \frametitle{Sommaire}
  \tableofcontents[currentsection]
\end{frame}

%\renewcommand\vec{\mathbf}
\renewcommand\vec{}

% ----------------------------------------------------------------------
\begin{frame}
  \frametitle{IA et régression}

  \begin{block}{Apprentissage automatique supervisé}

    \begin{itemize}
      \item données : $d+1$ variables $\vec{x} \in \R^d$ et $y \in \R$ observées sur $n$ individus. 
      \item phase d'apprentissage : déterminer les paramètres $(w_1, \dots, w_m)$ d'un modèle $f$
        qui calcule $y$ à partir de $\vec{x}$.
      \item phase d'inférence : étant donné une nouvelle valeur $\vec{x}^0$,
        calculer $f(w_1, \dots, w_m)(\vec{x}^0)$
        (alors que la valeur $y^0$ correspondante est inconnue). 
    \end{itemize}
  \end{block}

\end{frame}

% ----------------------------------------------------------------------
\begin{frame}
  \frametitle{Intuition dans $\R^1$ avec un modèle linéaire}

  \begin{center}
      \includegraphics[width=0.5\textwidth]{linear-regression}    
  \end{center}

  \begin{itemize}
  \item en bleu, les couples $(x_i,y_i)$ observées,
  \item en rouge, le graphe de la fonction $f(w_1, \dots, w_m)$
    qui donne $y$ à partir de $x$.
  \item $m = 2$ suffit dans ce cas (pente
    et ordonnée à l'origine).
  \end{itemize}
\end{frame}

% ----------------------------------------------------------------------
\begin{frame}
  \frametitle{Régression et optimisation}

  \begin{block}{Apprentissage $=$ optimisation}

    \begin{itemize}
    \item on considère une fonction \emph{de perte} qui
      est typiquement la moyenne des carrés des écarts (MSE)
      entre ce que donne le modèle et les observations :
      \[ L(w_1, \dots, w_m) :=
      \frac{1}{n} \sum_{i=1}^n (f(w_1, \dots, w_m)(\vec{x}_i) - y_i)^2 \]
    \item on la minimise par descente de gradient
      \begin{itemize}
      \item $f$ doit être différentiable et on doit calculer son gradient
      \item (outils de différentiation automatique, ex. \verb!PyTorch!)
      \item il y a des \emph{hyper-paramètres} ($\lambda$, nombre d'itérations, etc.)
      \end{itemize}
    \end{itemize}

  \end{block}
  
\end{frame}

% ----------------------------------------------------------------------
\newcommand\mat[1]{\mathbf{#1}}

\begin{frame}
  \frametitle{Exemple de modèle}
  
  \begin{block}{Définition sous forme matricielle}
    Soit $k \in \N$.
    On a $m = k(d+2)$ paramètres stockés dans les matrices :
    \begin{itemize}
      \item $\mat{M}_1$ de taille $k \times (d+1)$,
      \item $\mat{M}_2$ de taille $1 \times k$.
    \end{itemize}
    
    \[ f(\mat{M}_1, \mat{M}_2)(\mat{x}) := \mat{M}_2 \cdot \sigma ( \mat{M}_1 \cdot \mat{x} ). \]
    où $\sigma$ est une fonction d'activation non polynomiale appliquée
    élément par élément (ex. sigmoïde). 
  \end{block}
  
  \begin{block}{Théorème d'approximation universelle}
  Il existe $k$ tel qu'on approxime n'importe quelle fonction
  aussi bien qu'on veut par ce modèle ! 
  \end{block}

\end{frame}

%% TODO
%% % ----------------------------------------------------------------------
%% \begin{frame}
%%   \frametitle{Gradient ?}

%%   \begin{itemize}
%%   \item (immédiat ;-)) $\frac{\partial L}{\partial \mat{M}_2}(\mat{M}_1, \mat{M}_2, \mat{x})
%%     = \sigma ( \mat{x}^T \cdot \mat{M}_1^T )$
%%   \item (composition) on pose $\mat{z}:=\mat{M}_1\mat{x}$, $\mat{y}:=\sigma(\mat{z})$,
%%     \begin{align*}
%%     \frac{\partial L}{\partial \mat{M}_1}(\mat{M}_1, \mat{M}_2, \mat{x})
%%     &= \mat{M_2} \Bigg( \frac{\partial \mat{y}}{\partial \mat{z}} \frac{\partial \mat{z}}{\partial \mat{M}_1} \Bigg) \\
%%     &= \mat{M_2} \sigma'( \mat{M}_1 \cdot \mat{x} ) \mat{x}^T.
%%     \end{align*}
%%   \end{itemize}

%%   astuce : formulaire de dérivées matricielles.

%%   mieux : outils de dérivation automatique comme \verb!PyTorch!.
  
%% \end{frame}

% ----------------------------------------------------------------------
\begin{frame}
  \frametitle{Exemple de gros modèle : AlexNet (2012)}

    (CNN → RN → MP)² → (CNN³ → MP) → (FC → DO)² → Linear → softmax

où

\begin{itemize}
\item     CNN = convolutional layer (with ReLU activation)
\item     RN = local response normalization
\item     MP = maxpooling
\item     FC = fully connected layer (with ReLU activation)
\item     Linear = fully connected layer (without activation)
\item     DO = dropout
\end{itemize}

combinaisons de matrices et de fonctions d'activation \\
60 millions de paramètres 

\end{frame}
